\documentclass[a4paper,openright,9pt]{report}
\usepackage[utf8]{inputenc}
\usepackage[a4paper]{geometry}
\geometry{top=2cm, bottom=2cm, left=1.5cm, right=1.5cm}
\usepackage{graphicx}
\usepackage{booktabs}
\usepackage{wrapfig}
\usepackage[rightcaption]{sidecap}
\usepackage{float}
\usepackage{tabulary}
\usepackage{array}
\usepackage{multirow,array}
\usepackage{listings}
\usepackage{hyperref}
\usepackage{color} %red, green, blue, yellow, cyan, magenta, black, white
\usepackage{hyperref}
\usepackage{setspace}
\usepackage{caption}
\usepackage{subcaption}
\usepackage{placeins}
\usepackage[arrowmos]{circuitikz}
%otro paquete para circuitos que me parece mas facil
\usepackage{tikz}
\usepackage{fancybox}
\usepackage{array}
\usepackage{remreset}
\usepackage{titlesec}
\renewcommand{\thesection}{\arabic{section}}
\usepackage{comment}
\usepackage{siunitx}
\setcounter{secnumdepth}{4} %Para poner sub subsecciones (\subsubsection{<subsection>}) y para un nivel mas de subsección hay que poner \paragraph{<subsubsubsection>}
\titlelabel{\thetitle.\quad} %Para que ponga un punto desdpues de numerar una sección
\renewcommand{\figurename}{Figura} %En vez de que diga Figure aparece Figura, debajo de una figura
\usepackage{mathtools}
\DeclarePairedDelimiter{\abs}{\lvert}{\rvert} %definir modulo como \abs{}
\usepackage{physics} %Este paquete tiene derivadas mas chetas:
%	\dv{Q}{t} = \dv{s}{t}		 	para derivadas comun
%	\dv[n]{Q}{t} = \dv[n]{s}{t} 	para derivadas comun sucesivas (segunda derivada, tercera, etc..)
%	
%	\pdv{Q}{t} = \pdv{s}{t}  		lo mismo para derivadas parciales
%	\pdv[n]{Q}{t} = \pdv[n]{s}{t} 
%	\pdv{Q}{x}{t} = \pdv{s}{x}{t} 	para derivadas parciales multiples distintas
\newcommand{\code}[1]{\texttt{#1}} %Para escribir texto como codigo \code{<codigo>}

\renewcommand{\tablename}{Tabla} %Para que diga "Figura" en vez de "Figure"


\ctikzset{bipoles/voltmeter/text rotate/.initial=0,% <=new key
rotationvolt/.style={bipoles/voltmeter/text rotate=#1},% style for ease introduction in code
}

% code from pgfcircbipoles.sty
\makeatletter
\pgfcircdeclarebipole{}{\ctikzvalof{bipoles/voltmeter/height}}{voltmeter}{\ctikzvalof{bipoles/voltmeter/height}}{\ctikzvalof{bipoles/voltmeter/width}}{
    \def\pgf@circ@temp{right}
    \ifx\tikz@res@label@pos\pgf@circ@temp
        \pgf@circ@res@step=-1.2\pgf@circ@res@up
    \else
        \def\pgf@circ@temp{below}
        \ifx\tikz@res@label@pos\pgf@circ@temp
            \pgf@circ@res@step=-1.2\pgf@circ@res@up
        \else
            \pgf@circ@res@step=1.2\pgf@circ@res@up
        \fi
    \fi

    \pgfpathmoveto{\pgfpoint{\pgf@circ@res@left}{\pgf@circ@res@zero}}       
    \pgfpointorigin \pgf@circ@res@other =  \pgf@x  \advance \pgf@circ@res@other by -\pgf@circ@res@up
    \pgfpathlineto{\pgfpoint{\pgf@circ@res@other}{\pgf@circ@res@zero}}
    \pgfusepath{draw}

    \pgfsetlinewidth{\pgfkeysvalueof{/tikz/circuitikz/bipoles/thickness}\pgfstartlinewidth}

        \pgfscope
            \pgfpathcircle{\pgfpointorigin}{.9\pgf@circ@res@up}
            \pgfusepath{draw}       
        \endpgfscope    

    \pgftransformrotate{\ctikzvalof{bipoles/voltmeter/text rotate}}% <= magic line
    \pgfsetlinewidth{\pgfstartlinewidth}

    \pgfsetarrowsend{latex}
    \pgfpathmoveto{\pgfpoint{\pgf@circ@res@other}{\pgf@circ@res@down}}
    \pgfpathlineto{\pgfpoint{-\pgf@circ@res@other}{\pgf@circ@res@up}}
    \pgfusepath{draw}
    \pgfsetarrowsend{}
	

    \pgfpathmoveto{\pgfpoint{-\pgf@circ@res@other}{\pgf@circ@res@zero}}
    \pgfpathlineto{\pgfpoint{\pgf@circ@res@right}{\pgf@circ@res@zero}}
    \pgfusepath{draw}


    \pgfnode{circle}{center}{\textbf{V}}{}{}
}

\makeatother

\ctikzset{bipoles/ammeter/text rotate/.initial=0,% <=new key
rotationamper/.style={bipoles/ammeter/text rotate=#1},% style for ease introduction in code
}

\makeatletter
\pgfcircdeclarebipole{}{\ctikzvalof{bipoles/ammeter/height}}{ammeter}{\ctikzvalof{bipoles/ammeter/height}}{\ctikzvalof{bipoles/ammeter/width}}{
    \def\pgf@circ@temp{right}
    \ifx\tikz@res@label@pos\pgf@circ@temp
        \pgf@circ@res@step=-1.2\pgf@circ@res@up
    \else
        \def\pgf@circ@temp{below}
        \ifx\tikz@res@label@pos\pgf@circ@temp
            \pgf@circ@res@step=-1.2\pgf@circ@res@up
        \else
            \pgf@circ@res@step=1.2\pgf@circ@res@up
        \fi
    \fi

    \pgfpathmoveto{\pgfpoint{\pgf@circ@res@left}{\pgf@circ@res@zero}}       
    \pgfpointorigin \pgf@circ@res@other =  \pgf@x  \advance \pgf@circ@res@other by -\pgf@circ@res@up
    \pgfpathlineto{\pgfpoint{\pgf@circ@res@other}{\pgf@circ@res@zero}}
    \pgfusepath{draw}

    \pgfsetlinewidth{\pgfkeysvalueof{/tikz/circuitikz/bipoles/thickness}\pgfstartlinewidth}

        \pgfscope
            \pgfpathcircle{\pgfpointorigin}{.9\pgf@circ@res@up}
            \pgfusepath{draw}       
        \endpgfscope    

    \pgftransformrotate{\ctikzvalof{bipoles/ammeter/text rotate}}% <= magic line
    \pgfsetlinewidth{\pgfstartlinewidth}

    \pgfsetarrowsend{latex}
    \pgfpathmoveto{\pgfpoint{\pgf@circ@res@other}{\pgf@circ@res@down}}
    \pgfpathlineto{\pgfpoint{-\pgf@circ@res@other}{\pgf@circ@res@up}}
    \pgfusepath{draw}
    \pgfsetarrowsend{}


    \pgfpathmoveto{\pgfpoint{-\pgf@circ@res@other}{\pgf@circ@res@zero}}
    \pgfpathlineto{\pgfpoint{\pgf@circ@res@right}{\pgf@circ@res@zero}}
    \pgfusepath{draw}


    \pgfnode{circle}{center}{\textbf{A}}{}{}
}
\makeatother

\usepackage{etoolbox,refcount}
\usepackage{multicol}

\newcounter{countitems}
\newcounter{nextitemizecount}
\newcommand{\setupcountitems}{%
  \stepcounter{nextitemizecount}%
  \setcounter{countitems}{0}%
  \preto\item{\stepcounter{countitems}}%
}
\makeatletter
\newcommand{\computecountitems}{%
  \edef\@currentlabel{\number\c@countitems}%
  \label{countitems@\number\numexpr\value{nextitemizecount}-1\relax}%
}
\newcommand{\nextitemizecount}{%
  \getrefnumber{countitems@\number\c@nextitemizecount}%
}
\newcommand{\previtemizecount}{%
  \getrefnumber{countitems@\number\numexpr\value{nextitemizecount}-1\relax}%
}
\makeatother    
\newenvironment{AutoMultiColItemize}{%
\ifnumcomp{\nextitemizecount}{>}{3}{\begin{multicols}{2}}{}%
\setupcountitems\begin{itemize}}%
{\end{itemize}%
\unskip\computecountitems\ifnumcomp{\previtemizecount}{>}{3}{\end{multicols}}{}}


\begin{document}

\vspace{\fill}
\begin{figure}[H]
\centering
\includegraphics[width=.5\textwidth]{logofiuba.jpg}
\end{figure}

\begin{center}

\Huge{Trabajo Practico Control Adaptativo}
\\
\Huge{Modelado y Simulación de un Péndulo invertido}
\bigskip \\
\small{Gaggino, Lucas - Padrón: 99110}
\\

\small{\textit{86.20 Control Adaptativo }}
\\
\small{\textit{Departamento de Electrónica, Facultad de Ingeniería, UBA}}\\
\indent  

\end{center}
\vspace{\fill}
\newpage
\section{Modelado y simulación de la Planta }
\label{sec:red_hogar}

\subsection{Modelo y supuestos}
Consideramos un sistema plano compuesto por:
\begin{itemize}
  \item Un \textbf{carrito} de masa $M$ que se desplaza horizontalmente con coordenada $x$.
  \item Un \textbf{péndulo uniforme} (barra) de masa $m$ unido al carrito mediante una rótula ideal en su extremo superior. 
\end{itemize}
Definimos:
\begin{itemize}
  \item $\theta$ como el ángulo del péndulo \emph{medido desde la vertical hacia arriba}, de modo que $\theta=0$ es el equilibrio \textbf{inestable} (péndulo ``hacia arriba'').
  \item $l$ como la distancia desde el pivote al centro de masa de la barra (la barra completa tiene longitud $2l=L$).
  \item Gravedad $g>0$ hacia abajo. No hay rozamientos ni holguras.
  \item La única entrada es la \textbf{fuerza horizontal} $F$ aplicada al carrito.
\end{itemize}
A continuación se muestra una imagen del sistema descrito en la figura(\ref{fig:sist}).

\begin{figure}[h]
    \centering
    \includegraphics[width=0.4\textwidth]{Sistema.png}
    \caption{Sistema físico}
    \label{fig:sist}
\end{figure}

\subsection{Cinemática}
La posición del centro de masa del péndulo es
\begin{equation}
    X_p = x + l\sin\theta,
\qquad
Y_p = l\cos\theta .
\end{equation}

Sus velocidades:
\begin{equation}
\dot X_p = \dot x + l\cos\theta\,\dot\theta,
\qquad
\dot Y_p = -\,l\sin\theta\,\dot\theta.
\end{equation}
Por tanto, la rapidez del centro de masa cumple
\begin{equation}
\dot X_p^2+\dot Y_p^2
= \dot x^2 + 2l\cos\theta\,\dot x\,\dot\theta + l^2\dot\theta^2.
\end{equation}


\subsection{Energías}
\paragraph{Energía cinética.}
El péndulo aporta energía cinética por traslación (de su CM) y rotacional (alrededor de su CM). 
Para una barra uniforme de longitud $2l$, el momento de inercia alrededor del CM es
\begin{equation}
I_c=\frac{1}{12}m(2l)^2=\frac{1}{3}ml^2.
\end{equation}

Entonces,
\begin{equation}
T = \tfrac12 M\dot x^2
  + \tfrac12 m\big(\dot x^2 + 2l\cos\theta\,\dot x\,\dot\theta + l^2\dot\theta^2\big)
  + \tfrac12 I_c\,\dot\theta^2
= \tfrac12 (M+m)\dot x^2
+ m l \cos\theta\,\dot x\,\dot\theta
+ \tfrac{2}{3} m l^2 \dot\theta^2 .
\end{equation}

\paragraph{Energía potencial.}
Tomando el cero de energía potencial en $Y=0$ y eje $Y$ hacia arriba:
\begin{equation}
V = m g Y_p = m g l \cos\theta.
\end{equation}


\section*{Lagrangiano y ecuaciones de Euler--Lagrange}
El Lagrangiano es \(L = T - V\).
Sea $\bm q=(x,\theta)$, $\bm Q=(F,0)$ las fuerzas generalizadas no conservativas. 
Las ecuaciones de Euler--Lagrange son
\begin{equation}
\frac{d}{dt}\left(\frac{\partial L}{\partial \dot q_i}\right)
-\frac{\partial L}{\partial q_i}
= Q_i,\qquad i\in\{x,\theta\}.
\end{equation}

\paragraph{Ecuación para $x$.}
\begin{equation}
\frac{\partial T}{\partial \dot x}=(M+m)\dot x + m l \cos\theta\,\dot\theta,
\quad
\frac{d}{dt}\Big(\frac{\partial T}{\partial \dot x}\Big)
=(M+m)\ddot x - m l \sin\theta\,\dot\theta^2 + m l \cos\theta\,\ddot\theta,
\end{equation}
\begin{equation}
\frac{\partial L}{\partial x}=0,
\qquad
Q_x=F.
\end{equation}
Por tanto,
\begin{equation}
\boxed{(M+m)\ddot x + m l \cos\theta\,\ddot\theta - m l \sin\theta\,\dot\theta^2 = F.}
\label{eq:x}
\end{equation}

\paragraph{Ecuación para $\theta$.}
\begin{equation}
\frac{\partial T}{\partial \dot\theta}
= m l \cos\theta\,\dot x + \tfrac{4}{3} m l^2 \dot\theta,
\quad
\frac{d}{dt}\Big(\frac{\partial T}{\partial \dot\theta}\Big)
= - m l \sin\theta\,\dot\theta\,\dot x + m l \cos\theta\,\ddot x
  + \tfrac{4}{3} m l^2 \ddot\theta,
\end{equation}
\begin{equation}
\frac{\partial L}{\partial \theta}
= \frac{\partial T}{\partial \theta} - \frac{\partial V}{\partial \theta}
= \big(- m l \sin\theta\,\dot x\,\dot\theta\big) - \big(- m g l \sin\theta\big)
= - m l \sin\theta\,\dot x\,\dot\theta + m g l \sin\theta,
\end{equation}

\begin{equation}
Q_\theta = 0.
\end{equation}

Luego,
\begin{equation}
m l \cos\theta\,\ddot x + \tfrac{4}{3} m l^2 \ddot\theta - m g l \sin\theta = 0,
\end{equation}

o, dividiendo por $m l$,

\begin{equation}
\boxed{\cos\theta\,\ddot x + \tfrac{4}{3}\,l\,\ddot\theta - g \sin\theta = 0.}
\label{eq:theta}
\end{equation}


\subsection{Forma matricial}
Multiplicando \eqref{eq:theta} por $m l$ y agrupando:
\[
\underbrace{\begin{bmatrix}
M+m & m l \cos\theta \\
m l \cos\theta & \tfrac{4}{3} m l^2
\end{bmatrix}}_{\displaystyle \bm M(\bm q)}
\begin{bmatrix}\ddot x\\ \ddot\theta\end{bmatrix}
+
\underbrace{\begin{bmatrix}
- m l \sin\theta\,\dot\theta^2\\[2pt] 0
\end{bmatrix}}_{\displaystyle \bm C(\bm q,\dot{\bm q})}
+
\underbrace{\begin{bmatrix}
0\\[2pt] - m g l \sin\theta
\end{bmatrix}}_{\displaystyle \bm G(\bm q)}
=
\underbrace{\begin{bmatrix}
F\\ 0
\end{bmatrix}}_{\displaystyle \bm \tau}.
\]

Se observa que se obtiene la ecuación de Euler–Lagrange en forma canónica donde $M(q)$ es la matriz de inercia,   $C( q,\dot{ q})$ es la matriz de Coreolis donde se ven los efectos relacionados a la acelaración angular y si hubiera rozamiento viscoso también se vería en dicha ecuación. $G(q)$ Como el vector de fuerzas generalizadas  donde solo se ve la acción gravitatoria y $\tau$ el vector de fuerzas de control. 


\subsection{Aceleraciones explícitas (modelo no lineal)}
A partir de \eqref{eq:x}--\eqref{eq:theta}, resolvemos para $\ddot \theta$ y $\ddot x$.
Definamos
\[
\text{temp} = \\; \frac{F + m l \sin\theta\,\dot\theta^{2}}{M+m},
\qquad
\text{den} = l\left(\frac{4}{3} - \frac{m\cos^{2}\theta}{M+m}\right).
\]
Entonces
\[
\boxed{\;\ddot\theta \;=\; \frac{g\sin\theta - \cos\theta\,\text{temp}}{\text{den}}\;},
\qquad
\boxed{\;\ddot x \;=\; \text{temp} - \frac{m l \cos\theta}{M+m}\,\ddot\theta\;}.
\]
Obsérvese que, para pequeñas oscilaciones alrededor de $\theta=0$ y $F=0$, 
\(
\ddot\theta \approx \frac{g}{(4/3)l}\,\theta
\),
lo que confirma la \emph{inestabilidad} del equilibrio superior.

\subsection{Estado para integración numérica}
Si definimos el estado
\(
\bm y = [\,x,\ \dot x,\ \theta,\ \dot\theta\,]^\top,
\)
entonces el sistema de primer orden queda
\[
\dot{\bm y} =
\begin{bmatrix}
\dot x\\[2pt]
\ddot x(x,\dot x,\theta,\dot\theta;F)\\[2pt]
\dot\theta\\[2pt]
\ddot\theta(x,\dot x,\theta,\dot\theta;F)
\end{bmatrix},
\]
con $\ddot x$ y $\ddot\theta$ dados por las expresiones no lineales anteriores. 
Este es el sistema que se integra con \texttt{solve\_ivp} partiendo, por ejemplo, del equilibrio inestable
\(
\theta(0)=0,\ \dot\theta(0)=0
\)
y una fuerza constante $F=1$~N aplicada al carrito y el resultado se ve en la figura(\ref{fig:sim_nl_1}). 


\begin{figure}[H]
    \centering
    \includegraphics[width=0.7\textwidth]{sim_no_l.png}
    \caption{Simulación sistema físico no lineal }
    \label{fig:sim_nl_1}
\end{figure}

\begin{lstlisting}[language=Python]
def cartpole_ode(t, y):
    x, xdot, th, thdot = y
    s, c = math.sin(th), math.cos(th)

    # Auxiliar
    temp = (F_const + m*l*thdot*thdot*s) / (M + m)

    # Denominador con inercia de barra uniforme
    denom = l * (4.0/3.0 - (m * c * c) / (M + m))

    # Aceleraciones
    thddot = (g * s - c * temp) / denom
    xddot  = temp - (m * l * thddot * c) / (M + m)

    return [xdot, xddot, thdot, thddot]

# Condicion inicial: equilibrio inestable (pendulo arriba)
y0 = [0.0, 0.0, 0.0, 0.0]  # x, xdot, theta, thetadot

# Integracion
T = 5.0
t_eval = np.linspace(0.0, T, 1001)
sol = solve_ivp(cartpole_ode, [0.0, T], y0, t_eval=t_eval)
\end{lstlisting}


\section{Linealización, planta continua y discretización (ZOH)}

Partimos de las ecuaciones no lineales (barra uniforme, sin rozamientos)
\begin{align}
(M{+}m)\,\ddot x \;+\; m l \cos\theta\,\ddot\theta \;-\; m l \sin\theta\,\dot\theta^2 &= F,
\label{eq:nl_x}\\
m l \cos\theta\,\ddot x \;+\; \tfrac{4}{3} m l^2\,\ddot\theta \;-\; m g l \sin\theta &= 0.
\label{eq:nl_th}
\end{align}

\subsection{Linealización alrededor del equilibrio superior}
Linealizamos en torno a $\theta\simeq 0$, $\dot\theta\simeq 0$ usando $\sin\theta \!\approx\! \theta$, $\cos\theta \!\approx\! 1$ y despreciando $\,\dot\theta^2$:
\begin{align}
(M{+}m)\,\ddot x + m l\,\ddot\theta &= F,
\label{eq:lin_x}\\
m l\,\ddot x + \tfrac{4}{3} m l^2\,\ddot\theta &= m g l\,\theta.
\label{eq:lin_th}
\end{align}
Eliminando $\ddot x$ entre \eqref{eq:lin_x} y \eqref{eq:lin_th} se obtiene una ecuación escalar para $\theta$ de la forma
\begin{equation}
\ddot\theta \;=\; A_\theta\,\theta \;+\; B_\theta\,F,
\qquad
A_\theta \;=\; \frac{3 g (M{+}m)}{\,l(4M{+}m)\,}, 
\quad
B_\theta \;=\; -\,\frac{3}{\,l(4M{+}m)\,}.
\label{eq:thetadd_lin}
\end{equation}
Nótese que $A_\theta>0$, luego el equilibrio $\theta=0$ es \emph{inestable} en el modelo lineal.

\subsection{Planta continua $G(s)$ (entrada $F$, salida $\theta$)}
Aplicando transformada de Laplace a \eqref{eq:thetadd_lin} (condiciones iniciales nulas), queda
\begin{equation}
\frac{\Theta(s)}{F(s)} \;=\; G(s) 
\;=\; \frac{B_\theta}{\,s^2 - A_\theta\,}.
\label{eq:G_cont}
\end{equation}
En código (Python--Control), esto se representa como
\[
\verb|G_s = ctl.TransferFunction([B_theta], [1, 0, -A_theta])|.
\]
Reemplazando por los valores del sistema 
\begin{itemize}
    \item$ M = 1 \ Kg$
    \item $m = 0.1 \ Kg$
    \item $l = 0.5 m$
    \item $g = 9.81 \ \frac{m}{s^2}$
\end{itemize}

\begin{equation}
G(s) \;=\; \frac{-1.463}{s^{2} - 15.79}
\end{equation}

\subsection{Discretización por retención de orden cero (ZOH)}
Para un tiempo de muestreo $T_s$, la planta discreta $H(z)$ se obtiene mediante la discretización ZOH de \eqref{eq:G_cont}:
\begin{equation}
H(z) \;=\; \mathcal{ZOH}\!\left\{\,G(s)\,\right\}_{T_s}.
\end{equation}
En código:
\[
\verb|G_z_zoh = ctl.c2d(G_s, Ts, method='zoh')|,
\]
reemplazando queda la siguiente transferencia tomando un $Ts=0.02s$
\begin{equation}
H(z) \;=\; \frac{-0.0002928\,z \;-\; 0.0002928}{z^{2} - 2.006\,z + 1},
\qquad T_s = 0.02~\text{s}
\end{equation}
\subsection{Respuesta al escalón (continuo vs. discreto)}
Sea un horizonte $t_{\max}$ y un paso $T_s$. Las curvas usadas en las figuras se obtienen con:
\[
\verb|t_cont, y_cont = ctl.step_response(G_s, T=np.linspace(0, t_final, 1000))|,
\]
\[
\verb|t_disc, y_disc = ctl.step_response(H_z, T=np.arange(0, t_final, Ts))|.
\]
Como $G(s)$ es inestable (polo real en $+\sqrt{A_\theta} = 3.9738$), ambas respuestas divergen a diferencia del sistema no lineal, esto se debe a que al linealizar respecto al punto de equilibrio, estas ecuaciones se vuelven invalidas al momento que el sistema de aleja de dicho punto de trabajo. La discretización ZOH preserva la inestabilidad para $T_s$ suficientemente pequeño.
\begin{figure}[H]
    \centering
    \includegraphics[width=0.7\textwidth]{step_OL.png}
    \caption{Respuesta al escalón del sistema Continuo y Discreto a Lazo abierto }
    \label{fig:step_ol}
\end{figure}

\subsection{Diagramas de polos y ceros}

A continuación se resumen los polos y ceros de la planta continua $G(s)$ y de la planta discreta $H(z)$ (ZOH, periodo $T_s$):

\paragraph{Planta continua $G(s)$}
\begin{equation}
\text{Polos}\big(G(s)\big) = \{\, 3.97387812,\; -3.97387812 \,\}, 
\qquad
\text{Ceros}\big(G(s)\big) = \{\} 
\end{equation}
Como uno de los polos está en el semiplano derecho, $G(s)$ es inestable.

\paragraph{Planta discreta $H(z)$}
\begin{equation}
\text{Polos}\big(H(z)\big) = \{\, 1.08272127,\; 0.92359874 \,\}, 
\qquad
\text{Ceros}\big(H(z)\big) = \{\, -1 \,\}.
\end{equation}
Al existir un polo fuera del círculo unidad, $H(z)$ es inestable y preserva la inestabilidad del continuo. También notar que aparece un cero que en el sistema continuo no existe.

\begin{figure}[H]
    \centering
    \includegraphics[width=0.9\textwidth]{pz_mapOL.png}
    \caption{Diagrama de polos y ceros del sistema a lazo abierto }
    \label{fig:step_ol}
\end{figure}





\newpage
\section{Anexo}
\subsection{Referencias}
\begin{itemize}
    \item  \url{https://ctms.engin.umich.edu/CTMS/index.php?example=InvertedPendulum&section=SystemModeling}
\end{itemize}
\end{document}

